
Foundation model pretraining commonly uses static domain mixing ratios while largely ignoring when to present different data sources. Non-convex SGD optimization exhibits path dependence, suggesting that early training phases may influence final representational geometry.

This paper reports proof-of-concept implementation of diversity-ranked domain scheduling, which orders training domains from high to low diversity (measured via corpus-level vocabulary entropy, syntactic complexity, and semantic spread) using smooth Gaussian-weighted transitions.

The proof-of-concept confirms that diversity-ranked scheduling is implementable and testable at foundation model scale. Systematic diversity quantification for 6 Pile domains yields clear high-to-low ranking (0.92 to 0.35). The curriculum scheduler correctly implements Gaussian-weighted domain transitions with proper normalization and minimum weight constraints (verified via unit tests). The complete training pipeline executes without errors (22/22 unit tests pass, smoke test operational). The evaluation framework integrates with standard benchmarks (MMLU, Big-Bench, domain tasks).

However, performance improvement claims remain hypotheses pending full-scale validation. The smoke test (10 steps, single run) serves only to verify pipeline correctness, not to demonstrate convergence or statistical significance. The composite benchmark score of 0.2558 reflects an untrained model and provides no evidence for or against the performance hypothesis (≥2.0\% improvement at 1B scale, ≥0.5\% at 7B scale).

The proposed gradient geometry mechanism is similarly unverified. All four causal steps—diversity shapes gradient covariance, early geometry persists, late specialization preserves breadth, broad geometry enables robust continual learning—require full-scale multi-checkpoint experiments with Participation Ratio, CKA, and Fisher Information measurements.

Planned full-scale experiments (4 conditions × 2 scales × 5 seeds = 40 runs, estimated 6-8 weeks) will test whether diversity-ranked scheduling delivers performance improvements. If validated, this work establishes temporal data composition as a design principle for foundation model pretraining, complementing existing static mixture optimization methods. If performance improvements do not materialize, the approach nonetheless provides a structured framework for exploring curriculum learning at domain scale, with proof-of-concept validation confirming feasibility.

The question posed at the outset—does temporal order in which we present training domains matter as much as their proportions—can now be rigorously tested at scale. The answer will emerge from full-scale experiments currently planned.
